The consortium demonstrates a unique interdisciplinary synergy, combining technical innovation with deep cultural heritage expertise to ensure the project’s success. The collaboration is structured around three strategic pillars:
\begin{itemize}
  \item \textbf{Technology and Publishing Synergy:} 
  Arts et Métiers (CAPLAB) provides cutting-edge XR tools and 4D reality capture capabilities (motion, gesture, sound), which are essential for creating high-fidelity digital twins. Springer Nature ensures that these digital assets and research outcomes act as fuel for an innovative Open Science platform, reaching the widest possible academic and public audience through open access mechanisms.
  \item \textbf{Mediterranean Heritage Expertise:}
  Sapienza Università di Roma and the National and Kapodistrian University of Athens (NKUA) bridge the gap between classical Greco-Roman archaeology and modern digital documentation. Their expertise in urban conservation and climate risk assessment provides the necessary scientific grounding for the technological validation in high-stakes, outdoor environments.
  \item \textbf{Pan-European Regional Synergy:} 
  ELTE BTK (Central-Eastern Europe) contributes specialized AI-based processing and digital humanities expertise through its National Laboratory for Digital Heritage. Sorbonne Université (Atlantic region) coordinates the documentation of medieval fortifications and coastal erosion.
\end{itemize}

This interdisciplinary synergy ensures that the project outcomes are scientifically accurate, technologically advanced, and socially relevant.